% Neutral source-inspired NEJM research-article layout. XeLaTeX, two passes.
% No official logo, real paper text, authors, data, figures, DOI or download metadata.
\PassOptionsToPackage{no-math}{fontspec}
\documentclass[UTF8,fontset=none,scheme=plain]{ctexart}
\usepackage[paperwidth=567bp,paperheight=756bp,left=31.5bp,right=31.5bp,top=53bp,bottom=54bp,headheight=14bp,headsep=13bp,footskip=24bp]{geometry}
\usepackage{fontspec,tabularx,array,booktabs,multicol}
% Shared dependencies and neutral material. Compile with XeLaTeX.
% No publisher logos, paper text, data, figures, or download metadata are included.
\usepackage{amsmath,amssymb,graphicx,xcolor,tikz,fancyhdr,needspace,xurl,hyperref,balance,etoolbox}
% The legacy balance package omits the first split box's final baseline depth.
% Reserve that measured depth rather than hiding a CJK overfull-box warning.
\makeatletter
\patchcmd{\@BAlancecol}{\setbox\@leftcolumn=\vbox to \@tempdima}{\advance\@tempdima by \dp1\relax\setbox\@leftcolumn=\vbox to \@tempdima}{}{\PackageWarning{cns-template}{Balance depth patch unavailable; inspect final columns}}
\makeatother
\usetikzlibrary{arrows.meta,positioning}
\usepackage[protrusion=true,expansion=false]{microtype}
\IfFontExistsTF{Noto Serif CJK SC}{%
  \setCJKmainfont{Noto Serif CJK SC}%
}{%
  \setCJKmainfont[BoldFont=FandolSong-Bold.otf,ItalicFont=FandolKai-Regular.otf]{FandolSong-Regular.otf}%
}
\IfFontExistsTF{Noto Sans CJK SC}{%
  \setCJKsansfont{Noto Sans CJK SC}%
}{%
  \setCJKsansfont[BoldFont=FandolHei-Bold.otf]{FandolHei-Regular.otf}%
}
\setmainfont{TeX Gyre Termes}\setsansfont{TeX Gyre Heros}
\setmonofont{TeX Gyre Cursor}
\definecolor{JournalBlue}{HTML}{0097C9}
\definecolor{JournalRed}{HTML}{B83443}
\definecolor{ScienceRed}{HTML}{CD1021}
\definecolor{CiteBlue}{HTML}{0000EE}
\hypersetup{colorlinks=true,linkcolor=JournalBlue,urlcolor=JournalBlue,
  pdftitle={CNS论文汉化：中性版式示例},pdfauthor={示例模板；非期刊官方模板},
  pdfsubject={Neutral source-inspired Chinese layout demonstration; no real research claims}}
\urlstyle{same}\raggedbottom
\setlength{\parindent}{1em}\setlength{\parskip}{0pt}
\setlength{\emergencystretch}{2em}\tolerance=2200
\XeTeXlinebreaklocale "zh"\XeTeXlinebreakskip=0pt plus .1pt
\pagestyle{fancy}\fancyhf{}
\renewcommand{\headrulewidth}{0pt}\renewcommand{\footrulewidth}{0pt}
\newcommand{\PaperTitle}{中文题名占位：从组织结构到证据解释}
\newcommand{\SampleAuthors}{Author A\textsuperscript{1}, Author B\textsuperscript{2,*}}
\newcommand{\SampleAffiliations}{\textsuperscript{1}作者单位一（示例占位）。\quad
  \textsuperscript{2}作者单位二（示例占位）。\par
  *通信作者信息待填。DOI、收稿和出版日期均应从目标原文核实。}
\newcommand{\SampleAbstract}{本段是中性摘要占位文字，不陈述任何真实研究发现。它演示中文题名、作者信息、摘要和双栏正文之间的层级关系。实际使用时，应完整替换为目标论文的忠实中文译文，并保留研究对象、方法、结果和结论的限定条件。图注、参考文献、作者声明和来源信息均属于独立的覆盖项。模板中的图形仅说明排版结构，不表示生物学机制、实验结果或统计关系。}
\newcommand{\DemoCaption}{\textbf{图1. 中性证据链示意（占位图，无实验数据）。}
  \textbf{A}，输入节点表示需要核对的原文材料。
  \textbf{B}，处理节点表示逐段转换及复核。
  \textbf{C}，输出节点表示正文、图注和书目等独立交付部分。
  箭头仅表示本示例的阅读顺序，不代表因果关系、效应方向或数量变化。
  该图由模板自行绘制，所有节点均为文字占位；无样本量、误差线、统计检验或比例尺。
  实际论文图注应逐项保留原稿报告的样本量、单位、检验方法、限定及图形署名。}
\newcommand{\DemoVector}{%
 \resizebox{\linewidth}{!}{%
 \begin{tikzpicture}[font=\sffamily\small,node distance=18mm,
   box/.style={draw=black!70,rounded corners=2pt,fill=black!3,minimum width=26mm,minimum height=16mm,align=center},
   flow/.style={-{Stealth[length=2.5mm]},line width=.7pt}]
  \node[box] (a) {\textbf{A}\quad 输入\\原文材料};
  \node[box,right=of a] (b) {\textbf{B}\quad 处理\\转换与复核};
  \node[box,right=of b] (c) {\textbf{C}\quad 输出\\中文译稿};
  \draw[flow] (a)--(b);\draw[flow] (b)--(c);
  \node[below=11mm of a,align=center] {标题与摘要\\正文与图表};
  \node[below=11mm of b,align=center] {保留限定\\检查单位};
  \node[below=11mm of c,align=center] {可追溯引文\\明确原文信息};
  \node[below=29mm of b,font=\sffamily\footnotesize,text=black!65] {示意占位；不包含真实实验数据};
 \end{tikzpicture}}%
}
\newcommand{\SampleBibliography}{%
 \section*{参考文献（示例占位）}
 {\fontsize{7.6bp}{10bp}\selectfont\setlength{\parindent}{0pt}\interlinepenalty=10000
 \hypertarget{ref:1}{}1. Author A. (20XX). Example reference title. \textit{Journal placeholder}. 卷、页码与DOI待填。此条仅供演示引文跳转。\par\smallskip
 \hypertarget{ref:2}{}2. Author B. (20YY). Another example reference title. \textit{Journal placeholder}. 出版信息待填；不要把占位条目带入正式译稿。\par}
}
\newcommand{\SampleSourceInfo}{%
 \section*{原文信息}
 以下字段全部为示例占位：作者；英文原标题；期刊；年份；卷期；页码或文章号；DOI或来源链接。实际使用时必须逐项替换并核对。\par
 本文件是\JournalName 风格的非官方中文译稿样板，依据一次已核验文章版式进行结构化改编，不代表该期刊所有年代或文体的统一标准，也不保证像素级复刻。无真实论文内容或期刊标识图像。\par
}
% Deliberately neutral demo prose, not translated from a published paper.
\newcommand{\SampleOpening}{%
 \section*{引言：示例正文与引文}
 这是一段用于检查双栏阅读节奏的中性示例文字。段落中的术语、缩写和标点需要在实际项目中与源文保持一致。中文译稿应呈现原论文的完整论述，而不是把多段证据压缩成一段摘要。这里的引文只链接到本文件末尾明确标为占位的参考文献\SampleCiteOne。\par
 若原句仅表示相关、可能或提示，译文也应保留相同的证据强度。数字、单位、正负号、上下标和范围端点应逐项核对。原文存在内部矛盾时，应分别保留相应位置的陈述，再用简短译注说明；不能根据常识把不同的数字静默统一。\par
 \subsection*{内容层级与跨页衔接}
 此小标题演示二级层级。实际论文可能采用不同的小标题密度，模板不会替代对原文结构的识别。双栏文章需按左栏、右栏和跨页续句的顺序读取。若段落在页末中断，应在完成整句核对后接入下一页，避免把图中文字误当正文，或重复翻译同一段。\par
 图表具有独立的信息范围。正文引用到某幅图，并不意味着其图注可以省略。图注中的每个分面、比较对象、时间点、样本量及统计说明都应对应源文。英文图内标签可以保留，并在图旁或文末提供中文释义；实际采用哪种安排，需要依据图形尺寸和译文长度判断\SampleCiteTwo。\par
 \subsection*{来源与可维护数据}
 本段继续演示正文密度。可编辑文本和图注数据应与最终版面分开维护，以便修订时准确定位变化。若只调整版式，应能证明正文、图注和书目内容没有被意外改写。更换字体或纸张后，须重新检查断行、图注与图片的相对位置以及页脚可读性。\par
 文内引用与书目之间应有明确对应关系。这里的两条演示书目没有真实DOI，也没有虚构的实验信息。使用模板时可增删条目，但必须同时更新编号、链接目标和引用位置。页码只表示中文译稿自己的分页，不用于冒充原出版页码。\par
 \subsection*{结果叙述的演示占位}
 本节的文字仅用于呈现长段落在双栏页面中的流动，不报告真实结果。实际翻译应保留观察对象、实验条件和结局指标之间的关系。例如，原文分别描述体外实验与动物实验时，不能把两类结果合并为同一层面的结论；研究使用的对照、时间顺序和分析范围也必须保留。图表中的比较范围不应扩大为对整个人群或全部疾病阶段的判断。\par
 论述中的转折与限制同样属于正文。若研究先提出一种解释，再讨论另一种可能，译文应沿用这一推进顺序。作者明确写出的不确定性应能在中文中被识别；如果尚未确定某个现象的方向或机制，就不能改成已经建立的因果关系。对于术语近似但研究层级不同的概念，可在首次出现时保留缩写，后文保持统一。\par
 \subsection*{方法信息与可复核性}
 方法部分往往包含密集数字和专有名词。翻译时须区分浓度、总量、频率和处理时长，不能因为中文句式调整而改变参数的归属。试剂名称、软件版本和数据访问标识通常保留原样；上下限、排除规则和显著性标准则需要逐字符核对。这里没有提供任何真实实验参数，避免把模板占位误认为可以直接执行的方法。\par
 表格与补充材料的引用也需要保留原层级。正文对补充图的指向并不意味着模板中已经附有该补充图；若源文件只包含主文，应如实标注所交付的范围。图表编号、附录顺序和参考文献数量都应由目标文章决定，本示例中的结构不构成固定配额，也不要求所有论文采用同样的章节。\par
}
\newcommand{\SampleClosing}{%
 \section*{讨论：版式适配与范围}
 上方矢量图是模板原创的排版占位图，不来自任何论文。图形与完整演示图注被放在同一页，便于读者直接核对分面说明。真实项目应使用经过授权或用户提供的原图素材，并保留图注署名与许可信息。尺寸不足时，可适度缩放或调整栏宽，不以删除图注解决空间问题。\par
 模板仅记录一次已核验文章的可观察设计特征。即使文章出自同一期刊，不同年代、研究论文与综述、校正清样与正式版本也可能采用不同的页眉和引文形式。应以用户提供的实际PDF为准，不能仅凭期刊名称推断配色、字号或引用方式\SampleCiteOne。\par
 \subsection*{演示检查项}
 本样板用于展示可编译性、中文字形、标题层级、跨栏图注、页眉页脚和内部链接。图中没有实际数据，因此不应被当作科学插图引用。替换内容后，还需检查所有页的渲染结果，确认无截字、碰撞、异常空白和遗漏段落。\par
 \SampleBibliography
 \SampleSourceInfo
}


\newcommand{\JournalName}{The New England Journal of Medicine}
\definecolor{NEJMRed}{HTML}{ED1C24}\definecolor{NEJMText}{HTML}{231F20}\definecolor{NEJMGrey}{HTML}{77787B}
\definecolor{NEJMCream}{rgb}{1,.995193,.942840}\definecolor{NEJMCaption}{rgb}{.959609,.941207,.918853}
\color{NEJMText}\hypersetup{linkcolor=NEJMText,urlcolor=blue,pdftitle={NEJM风格中文研究论文中性样板},pdfauthor={原创排版演示；非期刊官方模板}}
\newcommand{\JournalHeading}{\textcolor{NEJMRed}{\resizebox{196bp}{!}{\rmfamily\scshape\fontsize{9bp}{11bp}\selectfont The new england journal of medicine}}}
\fancyhead[C]{\raisebox{8bp}{\ifodd\value{page}\rmfamily\fontsize{8.5bp}{10bp}\selectfont 中文研究论文样板\else\JournalHeading\fi}}
\fancypagestyle{opening}{\fancyhead{}\fancyhead[L]{\hspace*{100.5bp}\raisebox{8bp}{\JournalHeading}}}
\fancyfoot[C]{\sffamily\fontsize{6.3bp}{8bp}\selectfont n engl j med｜nejm.org｜出版信息待填}
\fancyfoot[L]{\sffamily\fontsize{6.3bp}{8bp}\selectfont 非官方中文样板}
\fancyfoot[R]{\sffamily\fontsize{6.3bp}{8bp}\selectfont 中文译页\thepage}
\renewcommand{\normalsize}{\fontsize{9.3bp}{12.4bp}\selectfont}\setlength{\parskip}{2bp}\setlength{\columnsep}{14bp}
\ctexset{section={format=\centering\sffamily\fontsize{12.2bp}{15bp}\selectfont,beforeskip=11bp,afterskip=5bp},subsection={format=\sffamily\bfseries\fontsize{9.1bp}{12bp}\selectfont,beforeskip=9bp,afterskip=3bp}}
\newcommand{\SampleCiteOne}{\textsuperscript{\hyperlink{ref:1}{1}}}\newcommand{\SampleCiteTwo}{\textsuperscript{\hyperlink{ref:2}{2}}}
\newcommand{\CaptionBox}[1]{\par\noindent\begingroup\setlength{\fboxsep}{7bp}\colorbox{NEJMCaption}{\begin{minipage}{\dimexpr\linewidth-14bp}\sffamily\fontsize{7.6bp}{9.2bp}\selectfont #1\par\end{minipage}}\endgroup\par}
\newcommand{\AbstractLabel}[1]{\par\vspace{6bp}{\sffamily\bfseries\fontsize{7.5bp}{9.5bp}\selectfont\textcolor{NEJMRed}{#1}}\par\vspace{2bp}}
\begin{document}\normalsize\thispagestyle{opening}\vspace*{-8bp}
\noindent\hspace*{30.5bp}\begin{minipage}[t]{336bp}
\setlength{\fboxsep}{2bp}\setlength{\fboxrule}{.4bp}
\noindent\fbox{\colorbox{NEJMCream}{\parbox[c][13.2bp][c]{\dimexpr\linewidth-8.8bp}{\centering\sffamily\fontsize{10bp}{12bp}\selectfont\textcolor{NEJMRed}{原始研究}}}}\par\vspace{7bp}
{\centering\fontsize{17.5bp}{21bp}\selectfont 中文研究题名占位：\par 从完整材料到可追溯的证据表达\par}\vspace{6bp}
{\centering\sffamily\fontsize{8.7bp}{10.6bp}\selectfont\color{NEJMGrey}Author A, M.D.\textsuperscript{1}；Author B, Ph.D.\textsuperscript{2}；\par Author C, M.D., Ph.D.\textsuperscript{1,*}\par}\vspace{8bp}
\noindent\colorbox{NEJMCream}{\parbox{\dimexpr\linewidth-4bp}{\centering\fontsize{8bp}{10bp}\selectfont 摘\quad 要}}\par
\end{minipage}\par
\noindent\hspace*{30.5bp}\begin{minipage}[t]{336bp}\vspace{0pt}\setlength{\parindent}{0pt}\setlength{\parskip}{0pt}
\fontsize{10bp}{13.2bp}\selectfont
\AbstractLabel{背景}
本文件是用于核验版式的原创中性样板，不报告真实研究结果。摘要首先说明研究问题和既有证据，再交代需要解决的具体缺口。实际汉化时，应完整保留原文对于研究对象、既往发现和知识边界的描述；不能以通用背景替换作者原有论证，也不从外部资料添加原稿没有报告的发现。题名、作者和单位均为占位，正式使用前须逐项替换。
\AbstractLabel{方法}
本段演示结构化摘要的方法栏目。真实译稿应交代研究设计、研究人群或实验对象、干预与比较、主要结局以及分析范围。原文中的样本量、剂量、时间点、单位、注册信息和统计方法均须忠实保留。本样板没有研究参与者、实验处理或统计检验，不提供可以据此实施的临床或实验方案。补充材料是否提供，应以本次实际收到的文件为准。
\AbstractLabel{结果}
本段仅检验中文结果段的字密度、标点和断行，不包含实验数据。正式稿须区分绝对变化和相对变化、组间比较和组内比较、观察到的趋势和经过检验的差异。应保留误差指标、区间端点、正负方向以及所有限制条件。图注和表下注释具有独立的信息范围，不能因正文已引用图表就省略；摘要中的数字也应与正文和对应表格分别核对。
\AbstractLabel{结论}
中性样板用于展示刊头红色、浅米色横栏、灰色作者信息、窄摘要区与右侧出版信息的关系。实际论文的结论应与作者的证据强度一致，保留可能性、局限和适用范围。此处不构成医学建议，不是期刊官方模板。仅供学术交流，出版物版权归原作者及出版商所有。
\end{minipage}\hfill\begin{minipage}[t]{127bp}\vspace{0pt}\sffamily\fontsize{7.5bp}{10bp}\selectfont
\textsuperscript{1}作者单位一（占位）；\textsuperscript{2}作者单位二（占位）。\par\medskip *通信作者地址与电子邮箱待填。\par\medskip 在线日期、DOI与版本信息按目标原文核实。\par\medskip 本示例无真实作者或研究资料。
\end{minipage}
\clearpage
\begin{multicols}{2}
\section*{引言}
本页正文为原创排版示例。研究论文通常先提出问题，说明已有证据能够回答什么，再确定尚待解决的部分。实际翻译应沿用原文的论述顺序，保留范围限定与证据层级。术语、缩写、基因或蛋白名称需要前后一致，文内引文则应链接到原有编号对应的条目\SampleCiteOne。

双栏页面按照左栏、右栏及跨页续句的顺序阅读。图中文字不属于邻近正文，下载来源标记也不应混入句子。若一句话跨页，应先确认首尾衔接，再决定中文分页位置；不能为凑版删除限定条件、压缩成摘要或重复接入同一段。
\section*{方法}
\subsection*{资料范围与组织}
本示例采用中性占位内容，不报告真实样本、干预或结果。真实项目应分别登记题名摘要、正文方法、完整图注、表格、书目、作者声明和原文出处。正文提到某份附录，不意味着该附录已随主文提供；仅有链接时，应在范围说明中明确。

可编辑文字与图像素材分别维护。图内保留英文时，应补充完整中文标签释义。复用原图前需要逐幅确定真实边界，保留全部面板、坐标轴和注释，并移除个人下载来源对象。仅修改裁剪框或覆盖白块不能证明隐藏标记已被删除\SampleCiteTwo。
\subsection*{阅读与数据核验}
数值核验关注单位、方向、范围和比较对象。表格中的百分比、总数与分组数应对应源文；原文已有疑点时，不得根据常识静默统一。排版调整完成后，应再次核对正文和书目的内容指纹，确认没有漏段、重复或字形丢失。
\end{multicols}
\vspace{5bp}\noindent\begin{minipage}{\linewidth}\centering\begin{minipage}{.82\linewidth}\DemoVector\end{minipage}\par\vspace{5bp}\CaptionBox{\DemoCaption}\end{minipage}
\begin{multicols}{2}
\subsection*{图文安排}
本页将一幅原创矢量占位图置于正文附近，完整图注采用浅米色底、无衬线小字号。实际译稿应根据图形尺寸和图注长度分配空间，避免把所有大图集中在文末。大图需要独立页时，应在正文中保持清楚的引用和顺序。
\end{multicols}
\clearpage
\noindent\textbf{\textcolor{NEJMRed}{表1．}中性内容覆盖表（占位，无研究数据）。}\par\smallskip
{\fontsize{8bp}{10.4bp}\selectfont\renewcommand{\arraystretch}{1.12}\noindent\begin{tabularx}{\linewidth}{@{}lXX@{}}\toprule
内容类别 & 核验对象 & 完成条件\\\midrule
正文 & 各节与跨页续句 & 保留完整顺序与限定\\
图表 & 全部分面、注释与单位 & 真实边界及可读性通过核验\\
书目 & 原编号与原始条目 & 引文能够跳转至对应条目\\
来源 & 原题、版本与出版信息 & 缺失字段明确，不编造\\\bottomrule
\end{tabularx}\par}
{\fontsize{7.6bp}{9.2bp}\selectfont 注：本表用于排版演示，所有条目均为原创占位，不表示实证结果。\par}
\begin{multicols}{2}
\section*{结果}
本段只展示表格后的正文衔接。真实结果应按照研究问题和测量对象组织，明确某个结论来自图中的哪一面板或表格的哪一行。对于不同实验层级、不同随访时点或不同分析人群的发现，应保持原文区分，不能合并为范围更广的结论。

表格具有自己的标题、变量名称与脚注。中文重排可以改变行高和列宽，以便在保留内容的前提下提高可读性；不能通过省略行列、缩短范围或改变单位来适应页面。图注中的误差线、统计检验及样本定义同样需要逐项核对。
\section*{讨论}
这一版式样板根据一次核验的研究论文外观提炼，不能代表该期刊所有年代、文章类型或出版版本。实际项目仍应量测题名区、栏目结构、配色与正文密度。中文字符宽度和译文长度会改变分页，所以应保留结构关系，并对所有输出页逐页检查。

正文、图、表与完整图注尽量交替呈现。减少空白时，应先调整段落流动和图表位置，再检查图注行距与小字号是否仍易读。不能用删去科学内容、裁掉面板或把原刊页码当作中文分页来解决空间问题。
\subsection*{局限与适用范围}
本示例没有真实研究资料，不提供疾病风险、疗效、安全性或因果结论。正式译稿必须保留作者明确写出的不确定性，并将未提供的补充文件与已经译出的主文清楚区分。
\section*{作者信息与声明}
\textsuperscript{1}作者单位一（占位）；\textsuperscript{2}作者单位二（占位）。通信、资助、利益冲突和数据共享信息均须按原文替换。本中性文件仅供学术交流；出版物版权归原作者及出版商所有。
\SampleBibliography
\SampleSourceInfo
\end{multicols}
\end{document}
